
\section{Background and Related Work}
\label{background}

\subsection{FPGA Accelerator Co-Design}

FPGA-based accelerators have become an increasingly important platform for deploying modern AI workloads due to their flexibility, reconfigurability, and ability to provide energy-efficient computation across diverse applications~\cite{jiang2025fpga}. However, developing custom FPGA accelerators remains a complex process involving hardware-software co-design, hardware verification, high-level synthesis (HLS), and iterative performance optimization~\cite{gibsonDLAS2025}.

To shorten this accelerator design process, the SECDA methodology~\cite{haris2021secda} was introduced to simplify accelerator development by enabling rapid hardware-software co-design using SystemC-based simulation and FPGA execution flows. 
The SECDA-TFLite~\cite{haris2023secda} toolkit further extended this ecosystem by enabling TensorFlow Lite-based AI accelerator development and deployment using reusable architectural templates and automated hardware generation flows. More recently, SECDA-LLM~\cite{haris2024designing} explored the integration of Large Language Models (LLMs) within the SECDA ecosystem to assist accelerator development workflows.

While these frameworks significantly reduce development overhead and improve accessibility of FPGA accelerator design, the exploration of accelerator configurations and hardware optimization strategies still largely depends on manual DSE, requiring expert knowledge to reason over architectural and hardware constraints.


%******************************************************************************************************

\subsection{DSE for FPGA Accelerators}

DSE is a widely used approach for identifying efficient hardware configurations under workload and device-specific constraints. Existing DSE methodologies for FPGAs commonly rely on exhaustive parameter sweeps, heuristic search methods, bayesian optimization, or techniques such as genetic algorithms to explore architectural parameter spaces~\cite{saeedi2024survey,biscontini2024machine}.

These approaches have proven their effectiveness in optimizing metrics such as latency, throughput, power consumption, and resource utilization. However, traditional DSE methods primarily operate over structured parameter spaces and often lack the ability to reason semantically about workload behavior, architectural patterns, or implementation trade-offs. Also, as the complexity of the accelerator increases, exhaustive exploration becomes computationally expensive due to the growing number of design permutations and hardware evaluation cycles~\cite{xu2025hardware}.

% SECDA-DSE aims to complement structured exploration approaches by integrating a reasoning-guided exploration layer via LLMs into the DSE workflow, enabling the system to leverage prior hardware knowledge and workload context when refining for accelerator configurations.
SECDA-DSE~\cite{sharma2026llm} provides an approach in which structured parameter exploration is augmented with an LLM-guided reasoning layer that can use prior hardware knowledge, retrieved design context, and evaluation feedback to refine accelerator configurations.


%******************************************************************************************************

\subsection{LLMs for Hardware and Systems Design}

%Recent advances in LLMs have demonstrated strong capabilities in code generation, structured reasoning, and software engineering tasks.

Recent advances in LLMs have motivated their use as assistants for engineering tasks that combine code generation, tool interaction, and structured decision making.
These capabilities have motivated growing interest in applying LLMs to hardware and system design workflows, including RTL generation, HLS assistance, compiler optimization, and architecture exploration~\cite{liao2024llms, li2025idse, zhang2026luminallmguidedgpuarchitecture}.

Several recent works have explored the use of LLMs for generating hardware descriptions or assisting accelerator development~\cite{fu2023gpt4aigchip, firouzi2024llm, vungarala2025sa}. Approaches such as RAG and CoT prompting also improve reasoning quality and contextual understanding for complex engineering tasks. In parallel, parameter-efficient fine-tuning techniques such as Low-Rank Adaptation (LoRA)~\cite{hu2022lora} enable adaptation of LLMs without requiring full model retraining.

Despite these advances, existing work largely focuses on isolated code generation tasks or standalone hardware synthesis assistance. Limited work has explored the integration of LLM-guided reasoning within structured FPGA DSE workflows that incorporate a retrieval-based context, iterative evaluation feedback, and real hardware execution.

In contrast, SECDA-DSE integrates structured design exploration, grounded retrieval reasoning, evaluation feedback, and hardware-aware refinement within a unified accelerator co-design framework. By combining SECDA-based hardware evaluation with LLM-guided reasoning, SECDA-DSE move towards adaptive and automated accelerator design workflows.